% SPDX-FileCopyrightText: Copyright (c) 2026 NVIDIA CORPORATION & AFFILIATES. All rights reserved.
% SPDX-License-Identifier: Apache-2.0
%
% ROLE: what the training and evaluation environment is, what it models and what it does not,
%   which workloads it replays at which loads, and how the cells were frozen. A reader should be
%   able to judge which results could fail to transfer (that judgment itself goes in sec:threats).
% CLAIMS: (1) AIS engine timing is in effect (setup check); (2) a replay is a deterministic
%   function of (policy, cell, replicate), so comparisons are paired on identical workloads;
%   (3) workloads, SLOs and loads were fixed on train segments with the default router only and
%   frozen, with SHA-256 hashes, before any policy was tuned.
% EVIDENCE: CR/config/engine.json; CR/facts/setup.json; CR/facts/noise.json;
%   CR/facts/noise_calibration.json; CR/facts/calibration.json; CR/facts/agentx_lowered.json;
%   CR/facts/test_freeze.json; CR/facts/remote.json; CR/facts/DEVIATIONS.md; CR/traces/MANIFEST.json;
%   CR/cells/{train,val,test}.jsonl and SPLIT_MANIFEST.json; CR/runs/calibrate-fix-r0/{curves,decisions}.json;
%   CR/audits/{setup,calibration,agentx-lowering}/; WT/benchmarks/learned_routing/learned_routing/
%   {replicates,cache,cells}.py and workloads/{synthetic,transform,agentx,agentx_lowered}.py;
%   WT/lib/mocker/src/replay/offline/README.md.
% STATUS: facts settled at lr-cells-v5 (frozen 2026-10-03 01:53 PDT). Statistics in tab:trace-stats
%   that the trace manifest does not record (pooled sessions, AgentX) were computed 2026-10-03 from
%   the trace files with the harness's own summarize(); each is marked in its src comment.
%
% Notation: only symbols from macros.tex / notation.tex. Concurrency, warm-up length, knee and
% think multiplier are written in words on purpose (no symbols for them in tab:notation).
% Restructure (2026-10-05): technical detail moved to appendix-workloads.tex (app:sim,
% app:workloads, app:calibration, app:freeze, sec:agentx); the AIS-league clause to a footnote
% (full text in appendix-ais-league.tex). sec:load-modes now labels setup.tex sec:loadmodes.

% The calibration figure, its colors, pgfplots styles and \simLevelLines moved to
% appendix-workloads.tex (app:calibration).

\section{Simulator and workloads}
\label{sec:simulator}

Every policy comparison in this paper, except the live check of \cref{sec:live}, comes from offline
replay in AISimulate, which provides both the replay engine (\aisim{}) and the forward-pass
performance model (\ais{}).
% src: WT/lib/mocker/src/replay/offline/README.md (replay loop and engines from aisimulate_core::replay);
%      CR/config/engine.json ais_perf_config_canonical_local (op-level timing); scope "except the live
%      check": restructure round 1 E1 (sec:live reports measured GPU runs)
\aisim{} runs Dynamo's KV router in front of $\Nw$ simulated vLLM engines, in simulated time, on recorded or generated request traces. This section describes that
environment and the workloads it replays, so that a reader can judge which results could change on
a live deployment. Three properties carry the rest of the paper. Engine timing comes from \ais{},
not from the mock engine's fallback model (\cref{sec:aisim}). A replay is a deterministic function of the policy, the cell
and a replicate index, so every comparison between two policies is paired on identical workloads
(\cref{sec:replicates}). And the cells, their service objective and their loads were fixed on
train segments with the default router alone, then frozen with SHA-256 hashes before any policy
was tuned (\cref{sec:calibration,sec:freeze}). What replay leaves out is listed in
\cref{sec:replay-gaps}.

% ---------------------------------------------------------------------------------------------
\subsection{Offline replay with \ais{} timing}
\label{sec:aisim}

% src: WT/lib/mocker/src/replay/offline/README.md ("The deterministic event loop, Generalized Mocker
%      Engine driver, logical-worker lifecycle, and report collector live in aisimulate_core::replay";
%      "virtual-time runtime"); CR/CONTRACT.md A9 addendum ("Offline replay runs in simulated time,
%      so host speed never changes results")
\aisim{}'s offline replay is a discrete-event simulation in virtual time
\citep{aisimulate,dynamo}. The event loop, the mock engines and the report collector
come from the \code{aisimulate-core} crate; Dynamo contributes the adapter that places each
arriving request through its KV router and feeds the engines' KV-cache events back into the
router's index. Simulated time jumps from one event to the next (an arrival, an engine step, a
completion), so the outcome of a replay does not depend on how fast the host runs it.
% src: CR/audits/setup/ais-and-policy-r0.md F4 (aisimulate-core src/engine/scheduler/vllm/core.rs);
%      CR/config/engine.json mock_engine_args (enable_chunked_prefill, enable_prefix_caching)
Each mock engine runs \code{aisimulate-core}'s model of the vLLM scheduler: continuous batching,
chunked prefill under a per-step token budget, and prefix caching over a paged KV cache. The
duration of each forward pass comes from \ais{}, an operator-level performance model of the
configured model, GPU, backend and parallelism.
% src: CR/facts/setup.json engine.ais_selected_estimation_mode = "op_level";
%      CR/config/engine.json ais_perf_config_canonical_local (database_mode SILICON, fallback_policy deny)
Batch-level simulators of this kind are an established tool for sizing LLM serving
\citep{agrawal2024vidur}.

% src: WT/notes/learned-routing/README.md Story; CR/CONTRACT.md "Policy YAML"
Worker selection in replay goes through the same Rust policy catalog and the same policy YAML as
the live router (\cref{sec:api-components}).
% Rejection of unknown types and parameters, the canonical YAML form and the main-branch clause
% moved to app:sim "Policies in replay" (main-branch clause also in background.tex's catalog footnote).

% src: CR/config/engine.json replay_call (top_level_ais_perf_config null, router_prefill_load_model
%      "none", note)
\ais{} times the engines only. The router-side \ais{} estimator, which some router modes use to
model prefill load, is off: the policies under study see the same router state a live router
without a performance model would see (\cref{sec:decision}).\footnote{A separate, secondary
  \ais{}-informed league of policies, which may read \ais{}-derived estimates at runtime, runs on
  its own build with the router-side estimator on (\cref{sec:excluded,sec:res-ais}).}
% src: CR/CONTRACT.md A16; prov: A16 (the secondary AIS-informed league; full text in
%      appendix-ais-league.tex, app:ais-league)

% src: CR/config/engine.json; CR/facts/setup.json (engine, kv_capacity)
\begin{table}[tbp]
  \centering\small
  \caption{The simulated deployment: engine configuration of every worker. Software versions and
    the configuration file's hash are in \cref{app:repro}.}
  \label{tab:engine}
  \begin{tabularx}{\linewidth}{@{}lY@{}}
    \toprule
    Model, backend & Qwen/Qwen3-32B on vLLM 0.24.0, aggregated (prefill and decode on one worker) \\
    Hardware & H100 SXM (\code{h100\_sxm}), tensor parallelism 2 \\
    Timing & \ais{}, operator-level estimation; router-side \ais{} off \\
    KV cache & block size 16; \num{18863} blocks (\num{301808} tokens) per worker, pinned to the
      \ais{} estimate so every replay uses the same capacity \\
    Context & \code{max\_model\_len} \num{131072} (YaRN) \\
    Batching & \code{max\_num\_seqs} \num{1024}; \code{max\_num\_batched\_tokens} \num{8192};
      chunked prefill and prefix caching on \\
    Workers & $\Nw \in \{2,4,6,8,16,32\}$ identical workers \\
    \bottomrule
  \end{tabularx}
\end{table}
% The Versions row (ai-dynamo-runtime 1.6.0; aisimulate 0.13.0.dev202609300000000061) and the
% "single source of truth: config/engine.json" caption note live in app:repro (Environment, Frozen
% inputs), restructure round 1 (E2).
% The TTFT comparison (AIS vs polynomial fallback) moved to tab:engine-ttft in appendix-workloads.tex.

\paragraph{Engine facts.}
% src: CR/facts/setup.json ais_active_evidence.estimator_vs_replay_ttft_ms
%      (1024: estimator 55.4868010845058, replay 55.486801; 8192: 475.97461213209505 vs 475.974612)
\Cref{tab:engine} lists the deployment. Setup confirmed that replay uses \ais{} timing: a
1,024-token request replayed alone matches the \ais{} estimate to within \SI{1}{ns}. The mock
engine's polynomial fallback underestimates a 120,000-token prefill 7.2 times
(\cref{tab:engine-ttft} in \cref{app:sim}), so replay without \ais{} would mislead about long-context traffic.
% src: CR/facts/setup.json single_request_ttft_ms_ais_vs_polynomial (120000: 17815.54 vs 2476.95;
%      ratio computed here: 7.19)
\Cref{app:sim} gives further engine facts.

% ---------------------------------------------------------------------------------------------
\subsection{What replay does not model}
\label{sec:replay-gaps}

Each item below is a difference between replay and a live deployment that could change which
policy wins; \cref{tab:caveats} in \cref{app:caveats} gives their status. \Cref{sec:threats} weighs them, and the
robustness and live checks of \cref{sec:robustness,sec:live} probe the first and the last directly.
\begin{itemize}
  \item \textbf{Router state is perfectly fresh.} Replay applies the engines' KV events to the
    router's index synchronously; a live router sees them late. Load counters come from the
    router's own request tracking in both settings, but whether their update timing matches is
    an assumption. A replay patch that delays these updates is used only as a test-time
    robustness check (\cref{sec:robustness}).
    % src: WT/lib/router-plugins/builtin/src/learned_choice/FEATURES.md "Information parity with the
    %      live router (LR-14)"; CR/facts/lag_patch.json semantics; CR/CONTRACT.md A5
  \item \textbf{One cache tier.} Replay models only the GPU (device) KV cache. A live router with
    host or disk tiers blends them into its overlap estimate, a signal no policy saw in training.
  \item \textbf{Fields a live router would not have.} Replay fills the request's expected output
    length with the true output length, and in closed-loop mode it passes the router a single-use
    session identifier for every flat trace row, which open-loop mode and a live client do not.
    No headline policy reads the output length (\cref{sec:integration}). A single-use session
    identifier never has a previous request, so session affinity is unaffected, and the \code{v2}
    home indicator keys on the prompt prefix only (\cref{app:policy}).
    % src: CR/facts/setup.json expected_output_tokens (populated_from_true_osl: true);
    %      CR/facts/UPSTREAM_FOLLOWUPS.md #6 (synthetic request_<line> session IDs in closed loop;
    %      campaign fix: v2 hash_home keyed on prompt prefix only; CR/facts/build_fix_r0.json)
  \item \textbf{\ais{} fidelity} [hypothesis about the size of the effect]. \ais{} times a
    prefill pass that mixes requests from the mean new and mean cached length of the batch, so
    adding a short prompt to a long-context chunk can shorten the predicted pass (from
    \SI{1744}{ms} to \SI{1150}{ms} in one probe). Its attention table for this configuration
    covers prompts up to 16,384 tokens with no cached prefix, so longer prefills are
    extrapolated. The likely effect is a small optimism when short prompts share a worker with
    long-context prefills; a search procedure could exploit it, which matters most for AgentX.
    % src: CR/audits/setup/ais-and-policy-r0.md F4 (7,936-token chunk at 100K prefix: 1,744 ms alone,
    %      1,150 ms with one 256-token prompt added; table covers isl <= 16,384 with prefix 0 only)
  \item \textbf{Errors near capacity.} Calibration places every cell near the default router's
    throughput knee (\cref{sec:calibration}), which is where routing matters and also where small
    timing errors are most amplified \citep{agrawal2024vidur}. This study does not validate
    \ais{}'s absolute accuracy; it relies on replay to rank policies, which the live check tests.
  \item \textbf{No network, no failures, identical workers,} and, in open-loop cells, arrivals
    that do not react to latency (\cref{sec:load-modes}).
\end{itemize}

% ---------------------------------------------------------------------------------------------
\subsection{Determinism and common-random-number replicates}
\label{sec:replicates}

\paragraph{Determinism.}
% src: CR/facts/setup.json determinism (per_request_identical for rr, lmetric, two-tier, seeded default,
%      AgentX; unseeded default 937-1995 rows differ) and replay_cost.in_process_repeats
%      (per_request_canonical_sha256 identical across 3 in-process repeats); determinism.decision
For a fixed policy, workload and policy seed, a replay reproduces every per-request row exactly,
across repeated runs in one process and across fresh processes. The only run-to-run variation
found at setup was the default picker's unseeded tie-break, which this study removes: the branch
adds a \code{seed} parameter to \defaultcf{}, and the harness always passes one. Bit-exact results
also need the same operating-system image (\cref{app:sim,app:parity}).

\paragraph{Where noise comes from.}
% src: CR/facts/noise.json why, finding_verdict; WT/benchmarks/learned_routing/learned_routing/
%      replicates.py (PROTOCOL, SPREAD_PROTOCOL, replicate_seed, policy_seed, permute_*,
%      spread_mooncake_arrivals); CR/CONTRACT.md A1
Because a replay is deterministic, noise has to come from the workload: what is arbitrary in a
recorded trace is the order of requests logged at the same instant. Replicate $\krep$ of a cell
re-draws only that arbitrary part (each session's first arrival inside its logging slot, or the
order of tied sessions or of recycled AgentX copies) together with the policy seed $\krep + 1$,
and the same replicate is used for every policy: common random numbers (\cref{app:sim}).
% The setup audit's rejection of run-twice noise estimates survives in tab:audits ("Setup,
% determinism, cost, context").
% src: CR/facts/DEVIATIONS.md 2026-10-02 calibration "arrival spread replicates" (N = 2, 4,000 input
%      tokens/s/worker: p85 slowdown 8.32 with bursts, 1.46 spread; ttft p50 1.5 s vs 92 ms;
%      runs/calibrate/sweep vs sweep2, mc-open N2)

Spreading arrivals inside their slot is a modeling choice, not only a noise source. Replay
delivers a logged burst at once whatever the arrival speedup, and with load matched per worker a
worker's share of each burst grows as $\Nw$ shrinks. With bursts kept, the default router's 85th
percentile slowdown at $\Nw = 2$ and 4,000 input tokens per second per worker was 8.32; with
arrivals spread it was 1.46. Without the spread, small-$\Nw$ cells would measure burst queueing
that the logging grid created.

Replicate noise is comparable to the effects of interest but smaller than segment heterogeneity
(\cref{sec:noise-floor,app:sim}). Replicates re-draw one segment; they do not create new
independent workloads (\cref{sec:eval}).

\paragraph{Cost.}
% src: CR/facts/setup.json replay_cost (fresh process 1.43-1.67 s for 2,000 requests at N = 2..32);
%      CR/facts/agentx_lowered.json caveats [8] (N = 32 lanes cells 2-4 min)
Replay is cheap enough for black-box search. A 2,000-request Mooncake replay takes about
\SI{1.5}{s} of one CPU core, nearly flat in $\Nw$; an AgentX cell at $\Nw = 32$ takes 2--4
minutes. \Cref{sec:compute} gives this study's compute.
% Full-trace times and per-request AgentX cost: app:sim "Cost". The cache key survives in
% appendix-reproducibility.tex app:determinism.

% ---------------------------------------------------------------------------------------------
\subsection{Workloads}
\label{sec:workloads}

The workloads differ in where prefix reuse comes from, because that is what a cache-aware router
exploits: shared system prompts across users (Mooncake, FAST25), a conversation's own growing
history (sessions), or long per-agent contexts carried across tool calls (AgentX).
\Cref{tab:trace-stats} summarizes the sources and \cref{tab:families} how they are split. A
\emph{segment} is a disjoint piece of one source, such as a time window, a generator seed or a
subset of agent plays. Segments are the unit of statistical independence: cells cut from one
segment at different $\Nw$, loads or transforms count as one workload (\cref{sec:eval}).

% src: CR/traces/MANIFEST.json files[*].stats (Mooncake, FAST25 rows: rows, isl, osl, first_arrival_span_ms,
%      mean_prefix_reuse_frac); sessions pooled over sessions_seed{0..9}.jsonl and AgentX over
%      traces/agentx/plays/*.json computed 2026-10-03 with learned_routing.workloads.common.summarize
%      (sessions: n 36,170, ISL mean 6,542.4 p50 5,539 p90 11,932, OSL mean 340.8 p50 249, reuse 0.830;
%       AgentX: n 3,132, ISL mean 66,313.0 p50 66,944 p90 103,488, OSL mean 748.7 p50 314,
%       within-play reuse 0.869 with transform.mooncake_stats' leading-block rule in source order)
\begin{table}[tbp]
  \centering\small
  \caption{Source traces as recorded or generated, before windows, transforms or recycling.
    Reuse is the mean fraction of a request's leading blocks already seen earlier in the trace,
    an infinite-cache hit rate; for AgentX it is computed within each play, since plays never
    share prefixes. Sessions are pooled over the ten \SI{720}{s} generator seeds.}
  \label{tab:trace-stats}
  \begin{tabular}{@{}lrrrrrrrr@{}}
    \toprule
    & & & \multicolumn{3}{c}{Prompt tokens} & \multicolumn{2}{c}{Output tokens} & \\
    \cmidrule(lr){4-6}\cmidrule(lr){7-8}
    Source & Requests & Span & mean & p50 & p90 & mean & p50 & Reuse \\
    \midrule
    Mooncake            & \num{23608} & 60 min & \num{8590}  & \num{6345}  & \num{16787}  & 182 & 30  & 0.63 \\
    FAST25 conversation & \num{12031} & 59 min & \num{12035} & \num{6909}  & \num{27367}  & 343 & 350 & 0.38 \\
    FAST25 synthetic    & \num{3993}  & 17 min & \num{15325} & \num{11587} & \num{38615}  & 149 & 69  & 0.42 \\
    Sessions (10 seeds) & \num{36170} & 12 min each & \num{6542} & \num{5539} & \num{11932} & 341 & 249 & 0.83 \\
    AgentX (82 plays)   & \num{3132}  & ---    & \num{66313} & \num{66944} & \num{103488} & 749 & 314 & 0.87 \\
    \bottomrule
  \end{tabular}
\end{table}

% src: CR/cells/{train,val,test}.jsonl (family, load.mode, segment counts verified 2026-10-03);
%      CR/cells/SPLIT_MANIFEST.json segments_by_split, agentx_play_subsets, segments
\begin{table}[tbp]
  \centering\small
  \caption{Segments and cells per split. Open: open-loop arrivals; closed: fixed concurrency;
    lanes: fixed number of concurrently running agent plays. FAST25 conversation shares Mooncake's
    arrival and length skeleton, so its windows are aligned to Mooncake's w3 and w5 and counted
    as those segments.}
  \label{tab:families}
  \begin{tabular}{@{}lllrrr@{}}
    \toprule
    & \multicolumn{1}{c}{Segments} & & \multicolumn{3}{c}{Cells} \\
    \cmidrule(lr){2-2}\cmidrule(l){4-6}
    Family & train / val / test & Load mode & Train & Val & Test \\
    \midrule
    Mooncake            & w0, w2, w4 / w1 / w3, w5 & open   & 12 & 3 & 18 \\
                        &                          & closed & 6  & 2 & 7  \\
    Sessions            & s0--s3 / s4, s5 / s6--s9 & open   & 6  & 3 & 9  \\
                        &                          & closed & 2  & 2 & 6  \\
    AgentX              & A1--A3 / V1--V3 / T1--T4 & lanes  & 8  & 4 & 14 \\
    FAST25 conversation & -- / -- / (w3, w5)       & open   & -- & -- & 2 \\
                        &                          & closed & -- & -- & 1 \\
    FAST25 synthetic    & -- / -- / x0, x1         & open   & -- & -- & 2 \\
                        &                          & closed & -- & -- & 1 \\
    \midrule
    All                 & 10 / 6 / 12              &        & 34 & 14 & 60 \\
    \bottomrule
  \end{tabular}
\end{table}

\paragraph{Mooncake.}
% src: CR/traces/MANIFEST.json mooncake (rows 23,608; block_size 512; distinct_roots 4;
%      first_arrival_span_ms 3,600,000; multi_turn_sessions 0); CR/cells/SPLIT_MANIFEST.json segments
%      mooncake:w0..w5; CR/facts/DEVIATIONS.md 2026-10-02 build fixer r0 (audit S1)
The public Mooncake trace \citep{mooncake} holds 23,608 single-turn requests over one
hour, with prefixes given as 512-token hash blocks under four distinct first blocks. It is cut
into six disjoint 10-minute slices, each a 4-minute warm-up followed by a 6-minute measurement
window: train w0, w2 and w4, validation w1, test w3 and w5.
% The earlier overlapping-slice design and its audit moved to app:workloads "Mooncake slices"
% (also recorded in tab:audits).

\paragraph{FAST25 conversation and synthetic.}
% src: CR/traces/MANIFEST.json roles and stats (fast25_conversation: skeleton shared with Mooncake;
%      fast25_synthetic: independent, max ISL 191,378); CR/cells/SPLIT_MANIFEST.json segments conv:w3,
%      conv:w5, fast25_synthetic:x0/x1
Two traces from the Mooncake project's FAST'25 release \citep{mooncake} are test-only families.
The conversation trace shares Mooncake's arrival grid and length skeleton but has a different
prefix structure; its windows are aligned to Mooncake's w3 and w5 and are counted as those
segments. The synthetic trace is independent of both, and its longest prompts exceed the context
limit and are dropped by the context-cap transform described below. A third trace from that release
is excluded as a relabeled copy of Mooncake (\cref{app:workloads}).

\paragraph{Synthetic sessions.}
% src: WT/benchmarks/learned_routing/learned_routing/workloads/synthetic.py (module docstring,
%      SessionSpec defaults); CR/traces/MANIFEST.json provenance.spec of sessions_seed{0..9};
%      CR/facts/DEVIATIONS.md 2026-10-02 build B2 "synthetic sessions are a generated Mooncake session trace"
Replay's built-in synthetic-session source cannot express a conversation: a later turn does not
extend the earlier turns' context, which leaves session affinity nothing to exploit. We therefore
generate multi-turn sessions in which every turn resends the whole conversation so far plus a new
user message, so its hash blocks extend the previous turn's. Each later turn is released a think
time after the previous turn \emph{completes}, so release is causal: a slow policy delays the
session's next request, as it would with a real user. One generator seed defines one segment
(s0--s3 train, s4 and s5 validation, s6--s9 test); \cref{app:workloads} gives the generator's
parameters.

\paragraph{AgentX.}
% src: CR/traces/MANIFEST.json agentx (source: semianalysisai/cc-traces-weka-062126 @ 23f152f6,
%      393 plays, 98,827 requests, 1,847,151,435 bytes; selection: predicate, 82 plays, 3,132 requests,
%      39 explicit subagent groups, no clipping or scaling); per-play stats.models (Claude models)
AgentX plays come from a public corpus of recorded agent sessions \citep{agentxtraces} (393 plays,
98,827 requests; \cref{app:workloads} gives the pinned revision). A play is one agent run: a dependency graph of requests, including subagents that
run in parallel and join, with recorded tool and think gaps between them. We keep the 82 plays
in which every request satisfies prompt plus $\max(\text{output}, 1) \le 131{,}072$, whole and
unchanged: 3,132 requests with 39 explicit subagent groups and a mean prompt of 66,313 tokens.
The plays were recorded with several Claude models and are replayed as Qwen3-32B traffic, so they supply
realistic structure (lengths, reuse, dependencies, gaps), not Qwen3-32B's behavior.
% src: CR/traces/MANIFEST.json agentx per-play stats.models over the 82 selected plays (7 distinct
%      Claude model IDs: Opus, Sonnet, Haiku and others; counted 2026-10-05)
Idle gaps over \SI{300}{s} are capped by a time warp that keeps every busy period and ordering
(\cref{app:workloads}).
% src: CR/cells/SPLIT_MANIFEST.json agentx_play_allocation ("seeded (agentx-split-v1) round-robin deal,
%      subagent plays first"), agentx_play_subsets (A1-A3 11 plays each; V1-V3 and T1-T4 7 each)
The plays are dealt, by a seeded round robin that places plays with subagents first, into ten
disjoint segments: A1--A3 (11 plays each) for training, V1--V3 and T1--T4 (7 each) for validation
and test. To keep 16 and 32 workers loaded, each play is converted once into \aisim{}'s agentic
trace format and recycled, with disjoint hash ranges so copies never share KV prefixes
(\cref{sec:agentx}).
% prov: A3 (AgentX recycling); the recycling subsection moved to appendix-workloads.tex.
% The load-mode subsubsection (sec:load-modes) merged into setup.tex sec:loadmodes, which now
% carries that label; its unique facts moved to app:workloads "Load-mode details".

\paragraph{Transforms.}
\label{sec:transforms}
% A run-in paragraph like the other workload paragraphs (restructure round 1, V9); the label now
% resolves to sec:workloads.
% src: WT/benchmarks/learned_routing/learned_routing/workloads/transform.py (module docstring, order 1-7);
%      CR/cells/SPLIT_MANIFEST.json train_ranges and transform_tags; CR/cells/test.jsonl transform cells
Seeded, deterministic transforms stretch prompt length (shared and unique parts separately), output
length, prefix-root count and think time, and cap context, so that the test can ask whether a
policy extrapolates (\cref{app:workloads}).
Train cells stay within $[1, 1.5]$ for both prompt multipliers, $[0.5, 2]$ for output and think
time, and $[1, 2]$ for prefix roots. Ten test cells go beyond these ranges: prompt multipliers of
2.5, output multipliers of 0.25, 2.5 and 4, four prefix roots, and think-time multipliers of
0.25 and 4.

% ---------------------------------------------------------------------------------------------
\subsection{SLO and load calibration}
\label{sec:calibration}

% src: CR/facts/calibration.json rules (levels, sla, windows), summary; CR/runs/calibrate-fix-r0/decisions.json
%      (binding decision table); CR/facts/DEVIATIONS.md 2026-10-02 calibration entries
Before any policy was tuned, calibration fixed each family's SLO pair $(\Imax,\Smax)$ and each
cell's load, using only $\piref$ (\defaultref) and round robin, and only train segments.
Calibration has to solve two problems at once. Routing only matters in a band of load: far
below capacity every policy meets the SLO, far above it every policy fails. And the SLO and the
band depend on each other, because whether a load is ``heavy'' depends on the SLO it is scored
against. The rules below fix both, and \cref{fig:load-curves} in \cref{app:calibration} shows the
result at $\Nw = 8$.

\paragraph{Load levels.}
% src: CR/facts/calibration.json rules.levels[0..3]
L1, L2 and L3 are the per-worker loads at which $\piref$'s windowed good fraction, averaged over
the train segments and two replicates, falls to 0.95, 0.85 and 0.65 under the family's SLO. Loads
are expressed per worker (offered input tokens per second for Mooncake-format open loop, sessions
per second for session open loop, concurrent sessions for closed loop, lanes for AgentX) and
multiplied by $\Nw$, rounding to an integer for closed loop and lanes. Each cell's measurement
window starts after a warm-up set per load mode (\cref{app:calibration}).

% src: CR/literature/LESSONS.md LR-12; CR/facts/calibration.json levels_per_worker
%      (mooncake open_speedup L2: N2 7476.77, N16 8993.04)
Matching load per worker does not hold the routing regime fixed across worker counts. Under
join-the-shortest-queue, waiting vanishes as the number of servers grows at a fixed per-server
load \citep{vanderboor2022survey}, and one measured learned router tied a simple baseline, led it
and tied it again as the pool widened \citep{tumkur2026calibrate}. Levels are therefore calibrated separately at every
$\Nw \in \{2,4,6,8,16,32\}$. They differ: Mooncake's open-loop L2, for example, is 7,477 input
tokens per second per worker at $\Nw = 2$ and 8,993 at $\Nw = 16$.\footnote{At $\Nw = 16$ and 32
  the AgentX grid starts at 3 lanes per worker, where $\piref$'s good fraction is already 0.916 and
  0.934, so AgentX has no L1 there and no cell uses one.}
% src: CR/runs/calibrate-fix-r0/curves.json agentx agentic_lanes by_n 16/32 (lowest grid load 3 lanes:
%      gf_default 0.916 and 0.934); CR/facts/calibration.json levels_per_worker agentx 16/32 L1 null
% src: CR/facts/calibration.json rules.levels[1] ("calibrated separately at each N ... on train
%      segments only"); CR/audits/calibration/independent-rerun-r1.md ("0 test touches in 30,060
%      calibration records"); DEVIATIONS 2026-10-02 calibration "levels are knee-matched per N"
Calibration therefore ran $\piref$ and round robin at the held-out worker counts, but only on
train segments, and it replayed no test trace.
% src: CR/facts/DEVIATIONS.md 2026-10-02 calibration "levels are knee-matched per N and per transform"
%      ("inheriting the untransformed level would have put them past the cliff"); calibration.json
%      rules.levels[2]. The 0.018 example DEVIATIONS cites there is also cited for the SLO headroom
%      (app:calibration), so it is used only once, there.
Transformed cells get their own levels, calibrated with the transform applied to the train
segments, because a transform changes how much load a worker can carry; the held-out FAST25
families inherit Mooncake's per-worker levels and SLO.

\paragraph{The SLO as a fixed point.}
% src: CR/facts/calibration.json rules.sla[0..5], sla.<family> (anchor_mode, fixed_point_history,
%      knee_fraction 0.9); CR/facts/DEVIATIONS.md 2026-10-02 calibration "SLA rule (A2 form, knee-anchored)"
Each family's bounds are set at $\Nw = 8$ in its primary mode: $\Smax$ so that $\piref$'s good
fraction is 0.65 at 0.9 times the family's load knee (where $\piref$ starts to saturate), which
makes that load L3, and $\Imax$ at 1.2 times the 95th percentile of $\piref$'s per-request mean
ITL at L2, iterated to a fixed point (\cref{app:calibration}). The resulting values are in
\cref{tab:calibration}.
% The ITL-alone failure rates at L2 (0.6--4.5%) survive in setup.tex sec:objective "Why two clauses".

% src: CR/facts/calibration.json summary.sla, summary.levels_n8_per_worker (rounded here);
%      CR/runs/calibrate-fix-r0/curves.json by_n["8"].levels.*.gf_rr_interp (Mooncake open 0.711,
%      sessions open 0.198, AgentX lanes 0.572)
\begin{table}[tbp]
  \centering\small
  \caption{Calibrated service objective per family and per-worker loads at $\Nw = 8$ (rounded).
    The SLO applies to every split, load mode and $\Nw$ of its family; the FAST25 families
    inherit Mooncake's. The last column is round robin's good fraction at $\piref$'s L2 in the
    primary mode, a measure of how much routing matters at that load.}
  \label{tab:calibration}
  \begin{tabular}{@{}lrrllrrrr@{}}
    \toprule
    & & & & & \multicolumn{3}{c}{Per-worker load at $\Nw = 8$} & RR good \\
    \cmidrule(lr){6-8}
    Family & $\Imax$ (ms) & $\Smax$ & Mode & Unit & L1 & L2 & L3 & frac.\ at L2 \\
    \midrule
    Mooncake & 252.117 & 3.46625 & open   & input tok/s & \num{5903} & \num{8798} & \num{11142} & 0.711 \\
             &         &         & closed & sessions    & 3.84  & 7.67  & 14.15 & \\
    Sessions & 46.0478 & 2.11784 & open   & sessions/s  & 0.350 & 0.380 & 0.409 & 0.198 \\
             &         &         & closed & sessions    & 29.0  & 33.1  & 37.5  & \\
    AgentX   & 38.6205 & 1.16134 & lanes  & lanes       & 2.57  & 3.39  & 4.74  & 0.572 \\
    \bottomrule
  \end{tabular}
\end{table}

\paragraph{What calibration does not fix.}
% src: CR/facts/calibration.json cache_pressure_open_loop.per_cell (55 open-loop
%      cells: min 0.616 sessions-s6-base-n32-open-L2, max 2.746 conv-w5-base-n8-open-L2; train
%      0.64-2.41, val 0.77-2.37, test 0.62-2.75) and lessons ("LR-12 action 2 ... no split has cells
%      below 0.5 (design gap recorded, not fixed)"; its 0.55 lower end predates the session fix)
Calibration fixes where the cells sit; it does not guarantee that every cell separates policies
(\cref{fig:load-curves}). One intended spread was not achieved. Measured as the distinct prompt
tokens touched in \SI{100}{s} of replay divided by the workers' total KV capacity, every one of
the 55 open-loop cells has a cache pressure between 0.62 and 2.75, so no split contains a cell
with little cache pressure.
% Moved to appendix-workloads.tex (app:calibration): log-load interpolation, knee rules, fixed-point
% and headroom history, fig:load-curves and its styles, the reading of the curves, realized bands,
% warm-up and measurement windows, the session warm-up fix, calibration compute (also in tab:cost).

% ---------------------------------------------------------------------------------------------
\subsection{Splits and the frozen test set}
\label{sec:freeze}

% src: CR/cells/SPLIT_MANIFEST.json holdout_axes, worker_counts, calibration.added_test_cells;
%      CR/cells/test.jsonl holdout_axis / holdout_axes / selection_exposed
The splits follow the segments of \cref{tab:families}. Training uses $\Nw \in \{4, 8\}$,
validation adds $\Nw = 6$, and the test set covers all six worker counts. Every test cell lies on a
test segment, so it is held out at least by workload; each also has a primary hold-out axis:
\begin{itemize}
  \item \textbf{time window:} Mooncake w3 and w5 at $\Nw \in \{4, 8\}$;
  \item \textbf{workload draw:} session seeds s6--s9;
  \item \textbf{agent plays:} AgentX segments T1--T4;
  \item \textbf{family:} FAST25 conversation and synthetic, never trained on;
  \item \textbf{transform extrapolation:} the ten cells beyond the train ranges
    (\cref{sec:transforms});
  \item \textbf{worker count:} $\Nw \in \{2, 6, 16, 32\}$, where $\Nw = 2$, 16 and 32 are
    extrapolation and $\Nw = 6$, which validation also uses, is labeled selection-exposed.
\end{itemize}
\Cref{tab:test-axes} in \cref{app:protocol} gives the number of test cells per axis and per $\Nw$. Calibration added four
AgentX test cells at $\Nw = 16$ and 32, which recycling made possible (\cref{sec:agentx}).

\paragraph{Freeze.}
% src: CR/facts/test_freeze.json (test_jsonl sha256, cells 60, traces, trace_sources,
%      cell_content_sha256, engine_json_sha256, split_manifest_sha256, bindings_build_id, rule);
%      CR/cells/SPLIT_MANIFEST.json calibration (cells_version lr-cells-v5);
%      CR/facts/test_results.json execution.evaluated_once ("149/149 hashes")
After calibration the cells were frozen. % prov: lr-cells-v5 (the frozen cell set)
The freeze recorded SHA-256 hashes of the test cells and their contents, the traces they replay,
the engine configuration and the split manifest. The freeze rule is that no test cell is evaluated
before the single test pass, which verified all 149 hashes before running. The test set was frozen
twice, once before and once after a calibration fix, with no test record in between
(\cref{app:freeze}).
% Hash prefixes, 39 files, 2.13 GB and 7 sources: in app:freeze and appendix-reproducibility.tex
% "Frozen inputs"; the first freeze's history: app:freeze.
